% Chapter 5 -- Rubric: PO4 (investigation, experiment, synthesis of conclusions),
% PO3 (3.4 evaluation against requirements), PO5 (application of modern tools)
\chapter{Investigation and Evaluation}
\label{ch:evaluation}

This chapter reports the investigative work carried out on the built system.
\Cref{sec:investigation} presents four investigations using analytical
derivation, structured adversarial inspection, automated test execution and
measured configuration evidence as appropriate to the question.
\Cref{sec:performance} evaluates the delivered system against every objective
and every functional and non-functional requirement of \Cref{ch:requirements}.
\Cref{sec:deviations} analyses each point at which the delivered system departs
from its specification and states the revision made or the reason none was made.

Pulse is a development-focused project rather than a research contribution, so
the investigations are validations of the design decisions that carried the most
risk rather than controlled experiments testing scientific hypotheses. Each
investigation states the form of evidence it can support; source inspection is
not presented as runtime measurement, and an unmeasured quantity is left
unresolved. One investigation (INV-2) tests the anonymity guarantee
adversarially, from the position of the most privileged insider.

Requirement identifiers (O-, FR-, NFR-, C-, G-, SH-, EC-) are those established
in \Cref{ch:problem,ch:research,ch:requirements}, so every claim here is
traceable back to a stated requirement. Repository-derived evidence was obtained
by executing the project's own tooling and inspecting the implementation paths
discussed below. The commands are listed in \Cref{subsec:reproduction} so the
figures can be reproduced. Where a quantity was not measured, the text says so
rather than estimating it.

\section{Experimental Investigation}
\label{sec:investigation}

\subsection{Selection of Investigations}
\label{subsec:investigation-selection}

Four decisions were judged to warrant investigation. A decision qualified if it
met at least two of three tests: its failure would invalidate a stated project
objective; its correctness is not observable from the user interface; and no
established practice prescribed the answer, so the team had to derive one.

\begin{longtable}{
  >{\raggedright\arraybackslash}p{0.07\textwidth}
  >{\raggedright\arraybackslash}p{0.33\textwidth}
  >{\raggedright\arraybackslash}p{0.17\textwidth}
  >{\raggedright\arraybackslash}p{0.30\textwidth}
}
\caption{Investigations and the design decisions they validate}
\label{tab:investigations}\\
\toprule
\textbf{ID} & \textbf{Decision under investigation} & \textbf{Objective / requirement at risk} & \textbf{Consequence if the decision is wrong} \\
\midrule
\endfirsthead
\toprule
\textbf{ID} & \textbf{Decision under investigation} & \textbf{Objective / requirement at risk} & \textbf{Consequence if the decision is wrong} \\
\midrule
\endhead
\bottomrule
\endlastfoot
INV-1 & Treating an absent metric by excluding it and renormalising the remaining
weights, rather than imputing a value for it & O-02, O-03, FR-22 & Projects with
partial tool coverage are scored unfairly, making the cross-project comparison
that is the product's central claim invalid \\
\addlinespace
INV-2 & Enforcing respondent anonymity in the database schema rather than in
application logic & O-04, NFR-07, NFR-08 & Respondents can be re-identified,
candour collapses, and the stakeholder conflict recorded in
\Cref{ch:requirements} (leadership insight vs.\ developer privacy) is
unresolved \\
\addlinespace
INV-3 & Using AI-assisted reranking to recover semantically related actions that
lexical matching can miss, while retaining lexical search as an independent
retrieval path & O-05, FR-39, NFR-13 & Without reranking, relevant actions
expressed with different vocabulary may be missed; when the provider is
unavailable, ordinary lexical retrieval remains available but explicit deep
search does not \\
\addlinespace
INV-4 & Bounded connector concurrency plus time-staggered scheduling to stay
inside free-tier API quotas & O-06, NFR-15, NFR-16 & A scheduled portfolio run
exhausts a shared credential's quota and the trend history that gives the product
its value develops gaps \\
\end{longtable}

\subsection{INV-1: Treatment of Absent Metrics in the Scoring Model}
\label{subsec:inv1}

\paragraph{The decision under investigation.} \Cref{ch:requirements} records
that most real projects have only a subset of the four tool categories
configured, and FR-22 requires that a score be computed from whichever metrics
are actually available. The scoring model must therefore define what a dimension
score \emph{means} when some of its input signals do not exist. Three treatments
were considered:

\begin{itemize}
  \item \textbf{Zero-imputation.} An absent signal scores 0 and retains its full
  weight.
  \item \textbf{Mean-imputation.} An absent signal takes the mean of the signals
  that are present, or of the portfolio.
  \item \textbf{Renormalisation.} An absent signal is excluded and the weights of
  the remaining signals are rescaled to sum to one. This is the treatment
  implemented by the shared scoring utility used across all risk dimensions.
\end{itemize}

The question is whether the choice is material, and if so which treatment is
defensible.

\paragraph{Why zero-imputation was rejected.} The objection is arithmetic rather
than empirical, and it can be shown with a single worked example using the
implemented weight vector of the security dimension
(\Cref{tab:security-weights}).

\begin{table}[H]
\centering
\caption{Implemented weight vector of the security dimension}
\label{tab:security-weights}
\begin{tabular}{llr}
\toprule
\textbf{Signal} & \textbf{Source tool} & \textbf{Weight} \\
\midrule
\texttt{securityRating} & SonarQube & 0.25 \\
\texttt{vulnCountDensity} & SonarQube & 0.20 \\
\texttt{securityReviewRating} & SonarQube & 0.15 \\
\texttt{securityHotspotsDensity} & SonarQube & 0.15 \\
\texttt{dependencyUpdateLag} & version control & 0.15 \\
\texttt{securityRemediationEffort} & SonarQube & 0.10 \\
\bottomrule
\end{tabular}
\end{table}

Consider a project that has connected version control but not SonarQube --- an
ordinary situation, and exactly the case FR-22 exists to handle. Five of the six
signals are absent, carrying 0.85 of the total weight. Suppose the one available
signal, \texttt{dependencyUpdateLag}, scores 80.

\begin{itemize}
  \item Under \textbf{renormalisation}, the single present signal's weight is
  rescaled from 0.15 to 1.0, and the dimension scores \textbf{80}. The score
  says: \emph{on the evidence available, this project's security posture looks
  good.}
  \item Under \textbf{zero-imputation}, the absent signals contribute
  $0.85 \times 0 = 0$ and the score is $0.15 \times 80 = 12$. The dimension
  scores \textbf{12}.
\end{itemize}

The same project, with identical underlying security practices, scores 80 or 12
depending only on whether an administrator has pasted in a SonarQube token.
Zero-imputation makes an absent measurement arithmetically indistinguishable
from a measurement of the worst possible value, so a project that has simply not
connected a tool is scored as though its code quality were \emph{known} to be
the worst observable.

This is disqualifying on two grounds. It violates FR-22 directly. More
seriously, it would defeat objective O-02: across a portfolio, the dominant
signal in the scores would be how many tools each project had configured, so the
cross-project comparison that the product exists to provide would be measuring
the wrong thing entirely. The distortion is also systematically
one-directional --- every absent signal drags the score down, never up --- so it
would not average out across a portfolio.

\paragraph{Why mean-imputation was rejected.} Mean-imputation avoids that bias
and, on accuracy alone, is a reasonable alternative to renormalisation. It was
rejected on a different ground.

Imputing the \emph{portfolio} mean makes one project's score depend on the
current state of every other project: a project's displayed score would change
when an unrelated project synchronised, without anything about the first project
having changed. That breaks O-03 and FR-26 outright, because a score breakdown
could no longer be stated in terms of that project's own metrics --- the honest
breakdown would have to read ``and 0.85 of this score is an assumption borrowed
from other projects.''

Imputing the project's \emph{own} mean avoids that coupling, but substitutes a
different problem: it fabricates evidence. It asserts that the project's
unmeasured security posture resembles its measured dependency hygiene, which is
an assumption the system has no basis for and cannot show the user.

Renormalisation makes the weaker and more honest claim --- \emph{this is the
score on the evidence available} --- and keeps every score a pure function of
that project's own inputs, which is what makes the per-signal breakdown of FR-26
possible at all.

\paragraph{Verification.} The implemented behaviour is asserted by the scoring
unit-test suite, which runs in CI on every change:

\begin{itemize}
  \item \texttt{renormalizedWeightedScore} \emph{``excludes missing signals and
  renormalizes the remaining weights to sum to 1''} (line 68) --- the direct
  assertion of FR-22.
  \item \emph{``computes a plain weighted average when every signal is
  present''} (line 56) --- confirms renormalisation is a no-op when nothing is
  missing, so the treatment does not perturb fully-configured projects.
  \item \emph{``returns null when no signal is present''} (line 48) --- a
  dimension with no evidence at all yields no score rather than a misleading
  zero, which is the same principle applied at the boundary.
  \item \emph{``excludes a present signal whose weight is 0''} and \emph{``clamps
  the result even if an individual signal score is out of 0..100''} (lines 80,
  89).
\end{itemize}

Each of the eight risk strategies additionally has a dedicated automated test
suite covering its signal set, satisfying NFR-22.

\paragraph{Limitation.} Renormalisation removes the bias, but it cannot
manufacture the information the absent signal would have carried. A dimension
score computed from two signals is a weaker estimate than one computed from six,
and the interface presents both identically --- a manager comparing two projects
cannot see that one score rests on materially less evidence. Team health is most
exposed, having only four signals, so a single absent input removes a quarter of
the evidence. This is recorded as an acknowledged limitation in
\Cref{subsec:deviation-confidence} and carried to future work.

\subsection{INV-2: Adversarial Investigation of Respondent Anonymity}
\label{subsec:inv2}

\paragraph{Research question.} NFR-07 requires that a survey response not be
attributable to the individual who submitted it. The absence of an identity
column establishes that no \emph{direct} attribution exists; it does not
establish that no attribution is \emph{possible}. The question is what an
adversary with realistic access could actually infer.

\paragraph{Threat model.} The adversary is the \textbf{company administrator}:
an insider with full read access to the tenant's data, including direct database
access, who wishes to learn how a named developer answered. This is the correct
adversary to model, because it is precisely the party whose access developers
(SH-03) are being asked to trust, and it is the strongest adversary the system
can meaningfully defend against. Defending against a database-level insider is
what distinguishes a structural guarantee from a policy promise.

\paragraph{Method.} Each plausible linkage route was enumerated and traced
through the applied response and recipient schemas, token representation,
dispatch path and insight-generation path. A route was classed as
\textbf{closed} only where the inspected representation contained no join key or
identity-bearing value; otherwise it was classed as \textbf{open} or
\textbf{statistically possible}. The result therefore establishes properties of
the evaluated implementation, not the probability of successful
re-identification in a real organisation.

\paragraph{Results.} \Cref{tab:anonymity} records each route attempted and its
outcome.

\begin{longtable}{
  >{\raggedright\arraybackslash}p{0.05\textwidth}
  >{\raggedright\arraybackslash}p{0.30\textwidth}
  >{\raggedright\arraybackslash}p{0.53\textwidth}
}
\caption{Re-identification routes attempted and outcomes}
\label{tab:anonymity}\\
\toprule
\textbf{\#} & \textbf{Linkage route} & \textbf{Outcome} \\
\midrule
\endfirsthead
\toprule
\textbf{\#} & \textbf{Linkage route} & \textbf{Outcome} \\
\midrule
\endhead
\bottomrule
\endlastfoot
A1 & \textbf{Direct.} A respondent column, foreign key or nullable author field
on \texttt{survey\_response} & \textbf{Closed by schema.} The table holds
\texttt{id}, \texttt{survey\_id}, \texttt{submission\_key} (a client-generated
UUID), \texttt{submitted\_at} and \texttt{answers} (jsonb). There is no identity
column and no foreign key to \texttt{User} \\
\addlinespace
A2 & \textbf{Token.} Recovering an individual from the access token embedded in
the survey link & \textbf{Closed by design.} \texttt{SurveyTokenPayload} is
\texttt{\{ surveyId, cycleId, deadline \}} only. One token is minted per
\emph{cycle}, not per recipient, so the token is byte-identical for every
developer on the project and carries no per-person information \\
\addlinespace
A3 & \textbf{Recipient join.} Correlating \texttt{survey\_recipient} rows
against \texttt{survey\_response} rows & \textbf{Closed.}
\texttt{survey\_recipient} records that a send was attempted and its outcome
(\texttt{sent}/\texttt{skipped}/\texttt{failed}). The only column it shares with
\texttt{survey\_response} is \texttt{survey\_id}, which is common to every
respondent, so the join is one-to-many in both directions and yields no
pairing \\
\addlinespace
A4 & \textbf{Timing.} Ordering responses by \texttt{submitted\_at} and
correlating against delivery times or individual working patterns &
\textbf{Open, statistical.} \texttt{submitted\_at} is stored at full precision
and is not coarsened. Because the link is broadcast once per cycle rather than
staggered per person, delivery time carries no per-person signal; but an
administrator could still correlate submission time against a developer's own
commit or CI activity, which the platform itself records. Narrowing rather than
identification, and harder as cohort size grows \\
\addlinespace
A5 & \textbf{Small cohort.} Inferring authorship on a project with very few
eligible respondents & \textbf{Open, statistical.} With few eligible respondents
the set of possible authors is small. This is a property of team size rather
than of anything the system stores, and narrows rather than identifies \\
\addlinespace
A6 & \textbf{Content.} A respondent identifying themselves within a free-text
answer, or writing in a recognisable style & \textbf{Open, irreducible.} No
technical control can prevent a respondent from identifying themselves in their
own words \\
\end{longtable}

\paragraph{Analysis and conclusion.} Route A1 is \textbf{closed structurally}:
direct attribution cannot be added without changing the response schema. Routes
A2 and A3 are closed by the current application and data design, but not by the
response schema alone. A future code change could mint a different token per
recipient, and a future schema or dispatch change could introduce a joinable
identifier. The defensible claim is therefore narrower: the stored response has
no direct identity field, while the shared-token and non-joinable-recipient
properties were verified by inspection of the evaluated implementation.

Routes A4--A6 are \textbf{statistical rather than structural}: they concern what
could be inferred from context, not what the system stores. \textbf{Respondent
anonymity is therefore enforced wherever the system can enforce it} --- no
stored response carries a respondent identity, the survey link is identical for
every recipient, and delivery records cannot be joined to responses. The
remaining routes depend on timing, team size and what respondents choose to
write, and no survey system can close them technically.

\subsection{INV-3: Availability of AI-Assisted Action Search}
\label{subsec:inv3}

\paragraph{Research question.} The action-search subsystem contains an ordinary
PostgreSQL lexical path and an explicitly requested \texttt{mode=deep} path that
reranks candidates through Pinecone. Lexical matching is fast and
provider-independent, but can miss conceptually related actions and problems
when they use different words; AI-assisted reranking exists to recover that
semantic similarity. A separate asynchronous pipeline generates Gemini
embeddings for actions. NFR-13 requires optional external-service
unavailability to produce graceful degradation. Does a deep-search request fall
back to lexical search when the reranker is unavailable, and what evidence
supports that behaviour?

\paragraph{Method.} The action-search endpoint was traced through its
request-handling and search-service layers. The acceptance criterion for full
conformance with NFR-13 was fixed before judging the result: when deep mode is
requested and Pinecone is unavailable, the endpoint must return lexical results
successfully and identify the serving path as lexical in its response metadata.
The three action-search test suites were then executed to determine whether that
failover is asserted automatically. Their test cases were reviewed as well as
counted, because a passing suite is evidence only for the behaviours it actually
asserts.

\paragraph{Results.} All 11 tests in the three suites pass (11/11, 836\,ms).
They assert query sanitisation and action scoping, Gemini request construction
and provider-response classification, and Pinecone result mapping and
malformed-score rejection. They do \textbf{not} invoke \texttt{searchActions()}
while the provider is unavailable, and therefore do not test lexical failover.

Inspection shows that the two paths currently remain separate rather than
switching automatically. With deep mode absent, the service calls PostgreSQL
lexical search and reports lexical mode. With deep mode requested, an
unavailable reranker produces an availability exception that the request layer
maps to HTTP \textbf{503} rather than initiating lexical retrieval. The
search-mode response metadata is set only after the service returns
successfully. Exposing that metadata cross-origin makes successful path
selection observable, but does not create automatic failover.

One availability property does hold independently: the embedding write path is
asynchronous, so logging an action does not block on Gemini. The embedding
record moves through pending, processing, ready and failed states, while a
content hash prevents unchanged actions from being embedded repeatedly. These
properties protect action creation, but they do not change the HTTP 503
behaviour of a requested deep search.

\paragraph{Analysis and conclusion.} NFR-13 is \textbf{partially satisfied}.
Ordinary lexical search remains usable without either AI provider, so the action
log remains searchable during an outage. A user who explicitly requests deep
search, however, receives a 503 when Pinecone is unavailable, so the experience
does not yet meet the stated graceful-degradation criterion. Full conformance
requires catching \texttt{ActionRerankingUnavailableError}, executing the
lexical query, returning HTTP 200 with \texttt{x-action-search-mode: lexical},
and adding integration tests for both unconfigured and temporarily unavailable
provider cases.

The honest limitation is that this establishes path availability, not retrieval
quality. The investigation does \textbf{not} show that Pinecone reranking
returns more relevant results than lexical search. Establishing that would
require a labelled query corpus with relevance judgements fixed in advance,
measured by recall@5 and mean reciprocal rank against the lexical baseline. That
corpus does not exist: the action log accumulates from first use (C-06) and
there is no partner action history to draw on.

A second exposure is that survey question generation and survey insight
generation also have no non-AI substitute: an extended Gemini outage postpones a
survey cycle rather than serving a reduced version of it. This is recorded as
residual risk R2 in \Cref{subsec:planning}.

\subsection{INV-4: Rate-Limit Headroom of Portfolio Synchronisation}
\label{subsec:inv4}

\paragraph{Research question.} Constraint C-03 confines the system to free
service tiers, and \Cref{ch:requirements} records the resulting conflict between
the client's expectation of current data and the providers' published rate
limits. Scheduled synchronisation (FR-18) fans out across every configured
project, and all projects imported from one workspace share a single access
credential. How many projects can one credential carry before a scheduled run
exhausts its quota?

\paragraph{Method and measured parameters.} The governing parameters were read
from the shipped configuration and code rather than assumed
(\Cref{tab:sync-parameters}).

\begin{longtable}{
  >{\raggedright\arraybackslash}p{0.30\textwidth}
  >{\raggedright\arraybackslash}p{0.28\textwidth}
  >{\raggedright\arraybackslash}p{0.30\textwidth}
}
\caption{Measured throttling and scheduling parameters}
\label{tab:sync-parameters}\\
\toprule
\textbf{Parameter} & \textbf{Value} & \textbf{Source} \\
\midrule
\endfirsthead
\toprule
\textbf{Parameter} & \textbf{Value} & \textbf{Source} \\
\midrule
\endhead
\bottomrule
\endlastfoot
Per-project tool fetch concurrency & 4 & Worker configuration, confirmed by
inspection \\
\addlinespace
Concurrency primitive & Bounded worker pool, rejects non-positive integers &
Shared synchronisation utility, unit-tested \\
\addlinespace
Spacing between projects sharing one credential & 10\,min default (bounded
0--240) & Scheduling configuration, confirmed by inspection \\
\addlinespace
Maximum scheduling spread & 180\,min default (bounded 0--1440) & Scheduling
configuration, confirmed by inspection \\
\addlinespace
Schedule time zone & Configurable, default UTC & Scheduling configuration,
confirmed by inspection \\
\addlinespace
Schedule times & Comma-separated HH:MM, validated and range-checked &
Configuration parser and validation tests \\
\addlinespace
Idempotency of a retried tick & Deterministic job ID per project and slot &
Synchronisation service and scheduled-worker tests \\
\addlinespace
GitHub GraphQL rate-limit guard & Pre-flight check of remaining quota against a
threshold & Version-control connector inspection \\
\end{longtable}

The scheduler assigns project index $i$ the delay $\min(iS, W)$. Where $S$ is
positive, this provides $\lfloor W/S \rfloor + 1$ distinct launch positions,
including time zero and the maximum-delay position. With the defaults
($S = 10$ minutes, $W = 180$ minutes), there are therefore \textbf{19 distinct
positions}: 0, 10, \ldots, 180 minutes. Projects 1--19 can occupy distinct
scheduled positions; project 20 and every subsequent project in the same
credential bucket are also assigned minute 180 and begin to overlap at the
scheduling level.

This calculation bounds launch spacing, not the number of simultaneously running
syncs. Actual overlap also depends on each sync's wall-clock duration, BullMQ
worker concurrency and queue backlog. Likewise,
\texttt{TOOL\_FETCH\_CONCURRENCY = 4} caps connector work \emph{inside one
project sync}; it is not a bound of four concurrent project syncs and cannot by
itself establish compliance with a provider quota.

\paragraph{Results and limitation.} The derivation above is sound but
\textbf{incomplete in one input}: the number of API calls issued per connector
per project sync was not instrumented, so the ceiling cannot be expressed in
projects-per-credential against a provider's published quota. Reading it off the
source is unreliable because the GitHub connector paginates and issues both REST
and GraphQL calls whose count depends on repository size.

\paragraph{Analysis and conclusion.} Two narrower conclusions are supported.
First, staggering reduces the initial request burst for the first 19 projects in
a credential bucket, while the per-project worker pool bounds simultaneous
connector fetches within each sync. Neither fact proves that rate-limit failure
becomes merely a longer run; that claim requires the missing call-count,
duration and worker-concurrency measurements. Second, the deterministic job
identifier makes a retried scheduled tick a no-op (NFR-12), so retrying the same
slot does not duplicate that slot's request volume. The projects-per-credential
ceiling remains unknown and is recorded as residual risk R1 in
\Cref{subsec:planning}.

\subsection{Synthesis}
\label{subsec:synthesis}

The four investigations support four bounded conclusions about the built system.
The scoring model's treatment of missing data is \textbf{sound by
construction} --- the rejected alternative would have scored an otherwise
identical project 12 instead of 80 on the worked example of \Cref{subsec:inv1},
purely for having one tool unconfigured --- which is the basis for the
cross-project comparability claimed by O-02. \textbf{Respondent anonymity is
enforced}: survey responses contain no identity field, the per-cycle token
carries no per-person information, and delivery records cannot be joined to
responses; the residual routes are statistical inferences from context rather
than stored linkages. AI-assisted reranking addresses semantic relationships
that lexical matching can miss; ordinary lexical search remains available
without AI, while explicit deep search currently returns 503 when Pinecone is
unavailable, so NFR-13 is partially met. Scheduled synchronisation has verified
staggering and idempotency controls, but its projects-per-credential capacity
remains unquantified until runtime request counts and durations are measured.

Collectively, the investigations also validate the decision recorded in
\Cref{ch:requirements} to prefer a rule-based scoring model over a learned or
LLM-judged one: INV-1 was possible \emph{at all} only because the model is
deterministic and decomposable, and neither of the rejected alternatives could
have been audited this way.

\section{Performance Evaluation}
\label{sec:performance}

\subsection{Setup, Tooling and Acceptance Criteria}
\label{subsec:eval-setup}

Every change reaching the main branch is validated by an automated pipeline that
is a required status check under branch protection. The evaluation below
therefore describes a branch state that is, by construction, the state of every
merged change rather than a specially prepared build.

The pipeline applies a path filter (\texttt{dorny/paths-filter}) so that a
backend-only change does not run frontend jobs and vice versa, with a
\texttt{concurrency} group that cancels superseded runs,
\texttt{permissions: contents: read}, and a single aggregate \texttt{ci-ok} job
that treats a skipped one-sided job as a pass and any failure or cancellation as
a failure.

For the backend the pipeline runs, in order: strict TypeScript type checking,
ESLint, a production build, \texttt{test:coverage} (Vitest with V8 coverage),
and \texttt{test:actions} (three \texttt{node:test} suites). For the frontend it
runs type checking, linting and a production build with
\nolinkurl{VITE_API_BASE_URL=/api/v1} baked in.

\paragraph{A note on the coverage figure.} The test runner is configured to
instrument \textbf{every backend application and library source unit}, rather
than only units imported by a test. This is a deliberate methodological choice
and it makes the headline number look far worse than a default configuration
would. Default instrumentation reports the percentage of \emph{the code the
tests already reach}, which flatters a codebase by silently excluding exactly
the modules nobody has tested. Instrumenting every source unit means an untested
module appears as 0\,\% rather than as an absence. The figure below is therefore
a percentage \textbf{of the system}, not a percentage of the tested part of it,
and \textbf{it is not comparable with a default-configured figure from another
project.}

Four verification methods are used in the tables that follow, each requirement
being assigned the weakest method sufficient to establish it:

\begin{itemize}
  \item \textbf{T --- Test.} An automated test asserts the behaviour and runs in
  CI on every change.
  \item \textbf{D --- Demonstration.} The behaviour was exercised end-to-end
  against the running system and observed.
  \item \textbf{I --- Inspection.} The property is established by reading the
  schema, configuration or source, and is not observable at run time.
  \item \textbf{A --- Analysis.} The property is established by derivation or by
  one of the investigations in \Cref{sec:investigation}.
\end{itemize}

Statuses are \textbf{Met}, \textbf{Partial} (met in part, or met by a weaker
mechanism than specified), \textbf{Not met}, and \textbf{Unverified} (not
established by any of the four methods; no claim is made).

\subsection{Measured Results}
\label{subsec:measurements}

All figures below were obtained by running the commands listed in
\Cref{subsec:reproduction}.

\begin{longtable}{
  >{\raggedright\arraybackslash}p{0.42\textwidth}
  >{\raggedright\arraybackslash}p{0.48\textwidth}
}
\caption{Quantitative results}
\label{tab:measurements}\\
\toprule
\textbf{Measure} & \textbf{Result} \\
\midrule
\endfirsthead
\toprule
\textbf{Measure} & \textbf{Result} \\
\midrule
\endhead
\bottomrule
\endlastfoot
TypeScript strict type check (backend) & \textbf{Pass} (exit 0) \\
ESLint (backend) & \textbf{Pass} (exit 0) --- no blocking findings and 131
warnings, predominantly \texttt{@typescript-eslint/no-explicit-any} \\
Vitest suites & \textbf{20 files, 131 tests, 131 passed, 0 failed}, 3.15\,s \\
\texttt{node:test} action-search suites & \textbf{11 tests, 11 passed, 0
failed}, 836\,ms \\
Total automated tests & \textbf{142 across 23 test files} \\
Statement coverage (all-files instrumented) & 12.37\,\% (722 / 5,835) \\
Branch coverage & 12.34\,\% (489 / 3,961) \\
Function coverage & 9.75\,\% (99 / 1,015) \\
Line coverage & 12.93\,\% (685 / 5,295) \\
Frontend automated tests & \textbf{None} \\
\end{longtable}

\begin{longtable}{
  >{\raggedright\arraybackslash}p{0.30\textwidth}
  >{\raggedright\arraybackslash}p{0.14\textwidth}
  >{\raggedright\arraybackslash}p{0.44\textwidth}
}
\caption{Statement coverage by area, showing where testing effort was
concentrated}
\label{tab:coverage-by-area}\\
\toprule
\textbf{Area} & \textbf{Statement coverage} & \textbf{Interpretation} \\
\midrule
\endfirsthead
\toprule
\textbf{Area} & \textbf{Statement coverage} & \textbf{Interpretation} \\
\midrule
\endhead
\bottomrule
\endlastfoot
Security and token utilities & \textbf{91.66\,\%} & The cryptographic core is
the most heavily tested code in the system, appropriately \\
\addlinespace
Version-control throttling and retry utility & 69.56\,\% & Provider throttling
and retry behaviour \\
\addlinespace
Risk-engine subsystem & 52.38\,\% & The eight strategies and scoring arithmetic
are covered; the dispatcher and descriptive signal catalogue are untested \\
\addlinespace
Synchronisation utilities & 52.17\,\% & Bounded-concurrency behaviour is
covered; connector registration is untested \\
\addlinespace
\texttt{GithubActions} / \texttt{Github} connectors & 14.67\,\% / 12.79\,\% &
Exercised only incidentally \\
\addlinespace
\texttt{Jira}, \texttt{SonarQube}, \texttt{GitLab}, \texttt{Linear} connectors &
0\,\% & Not unit-tested; validated by demonstration against live APIs instead \\
\addlinespace
Queueing, notifications, embeddings, authentication and reranking & 0\,\% &
I/O-boundary modules, exercised through doubles or by demonstration \\
\end{longtable}

\paragraph{Interpretation.} The distribution is more informative than the
headline. Testing effort went where correctness is both critical and cheaply
assertable --- cryptography, scoring arithmetic, concurrency bounds --- and the
near-zero areas are overwhelmingly I/O boundaries against third-party services,
which unit tests can only assert against a double anyway. That is a defensible
allocation under constraint C-02 (no dedicated QA member). It is nonetheless a
real limitation, and it should be stated plainly rather than explained away:
\textbf{NFR-22 requires scoring rules to be documented and covered by automated
tests and that is satisfied, but the system as a whole rests on 142 tests and
has no frontend test at all.} A regression in the React application would be
caught only by type checking, linting and manual use.

\paragraph{Security review.} \Cref{ch:requirements} commits the project to
OWASP-aligned practice. The team reviewed the implementation manually against
the OWASP Top~10 --- authentication and session handling, access control,
injection, cryptographic storage of connector credentials, and rate limiting of
the unauthenticated survey endpoint. The codebase was additionally analysed with
SonarQube through Pulse's own SonarQube connector, so the platform's security
dimension was applied to its own repository. This is an internal review rather
than an independent third-party assessment.

\paragraph{Partner acceptance.} The system was not piloted against SSCL's own
GitHub, Jira or SonarCloud data. Evaluation used the team's own repositories and
tool accounts, which exercised every connector end to end without bringing
partner data into the development environment. The system was instead reviewed
with SSCL's Chief Technical Officer in five meetings, each held after a major
improvement (\Cref{subsec:multidisciplinary}).

\subsection{Evaluation Against Functional Requirements}
\label{subsec:eval-fr}

\begin{longtable}{
  >{\raggedright\arraybackslash}p{0.07\textwidth}
  >{\raggedright\arraybackslash}p{0.06\textwidth}
  >{\raggedright\arraybackslash}p{0.57\textwidth}
  >{\raggedright\arraybackslash}p{0.13\textwidth}
}
\caption{Verification of functional requirements}
\label{tab:fr-verification}\\
\toprule
\textbf{ID} & \textbf{M.} & \textbf{Evidence} & \textbf{Status} \\
\midrule
\endfirsthead
\toprule
\textbf{ID} & \textbf{M.} & \textbf{Evidence} & \textbf{Status} \\
\midrule
\endhead
\bottomrule
\endlastfoot
FR-01 & I, D & Registration creates a company and administrator account; a
one-time verification code is retained transiently and the account is created
only after verification & Met \\
FR-02 & T, I & Login service and persisted server-side session store & Met \\
FR-03 & I & Single-use reset token; the reset flow explicitly invalidates every
existing session belonging to the user & Met \\
FR-04 & I, D & Single-use invitation accepted during registration or by an
authenticated account & Met \\
FR-05 & T, I & Central authorisation service and requester-role resolution;
member access gated on project-membership records & Met \\
FR-06 & D & Administrator removes a \texttt{projectmember} row & Met \\
FR-07 & D & Workspace endpoint binds a workspace to its provider, organisation
identifier and access token & Met \\
FR-08 & D & Workspace interface lists reachable repositories through the
provider adapter before any persistence & Met \\
FR-09 & D & Selected repositories imported as \texttt{project} rows & Met \\
FR-10 & D & Subsequent import excludes already-imported repositories & Met \\
FR-11 & D & Per-project \texttt{projecttoolintegration} rows for PM, quality and
CI/CD; VCS credential inherited from workspace & Met \\
FR-12 & I, D & Credentials masked by default; reveal is a separate
administrator-restricted action & Met \\
FR-13 & D & \texttt{POST /sync} enqueues a job and returns job ID and stream key
immediately & Met \\
FR-14 & T, D & Implemented connectors cover VCS, issue tracking, code analysis
and CI/CD & Met \\
FR-15 & I & One \texttt{projectsnapshot} row per run; raw metrics as JSONB in
four per-domain tables keyed by \texttt{snapshot\_id}; no update path & Met \\
FR-16 & D & SSE over \texttt{/progress}; per-tool status rendered by
\texttt{SyncProgressBanner} & Met \\
FR-17 & I, D & Persistent event store replays missed progress events to a
reconnecting client & Met \\
FR-18 & T, I & Schedule queue is reconciled from configured times when the
worker starts; covered by scheduled-worker tests & Met \\
FR-19 & I & Synchronisation processing tracks unsuccessful tools, marks each
tool's outcome, persists successful results and reports the run as partial &
Met \\
FR-20 & T & Eight strategy classes, each with its own automated test suite &
Met \\
FR-21 & T & Overall score composed from dimension scores per snapshot;
\texttt{riskscore.overall\_score} & Met \\
FR-22 & T, A & Scoring tests assert that missing signals are excluded and
remaining weights are renormalised to sum to one; design rationale in
\Cref{subsec:inv1} & Met \\
FR-23 & I & \texttt{riskscore} rows keyed to \texttt{project\_snapshot\_id}
(unique) & Met \\
FR-24 & D & Portfolio dashboard lists accessible projects with overall score,
dimension scores and trend & Met \\
FR-25 & D & Per-project view with current score, trend and dimension history
(Recharts) & Met \\
FR-26 & D, I & Score-breakdown interface renders contributing signals, values,
weights and contributions from the signal catalogue & Met \\
FR-27 & I, D & Survey-question service generates from health scores, trend and
notable metric conditions & Met \\
FR-28 & T & Deterministic deduplication followed by a scoring prompt with a
configurable minimum score of 60 & Met \\
FR-29 & I, D & Questions editable while \texttt{draft}; editing rejected once
\texttt{sent\_at} is set & Met \\
FR-30 & T, I & Distribution processor assigns a randomised monthly send moment
per project within a configurable window, fixed once assigned & Met \\
FR-31 & D & On-demand dispatch available to an authorised user & Met (pending
confirmation) \\
FR-32 & I, D & One anonymous link per cycle, broadcast to configured channels and
emailed individually to eligible developers & Met \\
FR-33 & I & Cooldown enforced from \texttt{survey\_recipient} history against
\texttt{SURVEY\_MIN\_DAYS\_BETWEEN\_SURVEYS} (default 15, bounded 1--60);
skipped sends recorded with \texttt{skip\_reason: 'cooldown\_active'}. Fair
rotation between a developer's projects is unconfirmed & Partial \\
FR-34 & T & Automated token tests cover decoding, cycle validation and deadline
expiry & Met \\
FR-35 & I, D & Survey-insight service produces category scores, themes,
narrative summary and per-question summaries on close & Met \\
FR-36 & I, D & \texttt{project\_survey\_status.pending\_survey} raised when
configured thresholds are breached & Met \\
FR-37 & D & \texttt{actions} row carrying problem, reason, action taken, date
and affected \texttt{project\_ids} & Met \\
FR-38 & D & Per-project and company-wide listings, scoped by role & Met \\
FR-39 & T, A & Lexical search and optional Pinecone reranking are implemented.
Ordinary lexical retrieval remains available without AI, but a requested deep
search returns 503 rather than falling back when Pinecone is unavailable;
investigation INV-3 (\Cref{subsec:inv3}) & Partial \\
FR-40 & I, D & \texttt{next\_review\_at} with deferral presets; non-modal
\texttt{WeeklyReviewBanner} and \texttt{InlineRating} & Met \\
FR-41 & I & Author may modify or delete own actions; administrator may do so for
any action in the company (\texttt{logged\_by\_user\_id},
\texttt{company\_id} on \texttt{actions}) & Met \\
\end{longtable}

\subsection{Evaluation Against Non-Functional Requirements}
\label{subsec:eval-nfr}

\begin{longtable}{
  >{\raggedright\arraybackslash}p{0.08\textwidth}
  >{\raggedright\arraybackslash}p{0.06\textwidth}
  >{\raggedright\arraybackslash}p{0.56\textwidth}
  >{\raggedright\arraybackslash}p{0.13\textwidth}
}
\caption{Verification of non-functional requirements}
\label{tab:nfr-verification}\\
\toprule
\textbf{ID} & \textbf{M.} & \textbf{Evidence} & \textbf{Status} \\
\midrule
\endfirsthead
\toprule
\textbf{ID} & \textbf{M.} & \textbf{Evidence} & \textbf{Status} \\
\midrule
\endhead
\bottomrule
\endlastfoot
NFR-01 & T, I & Credential-security subsystem encrypts connector secrets at rest
with AES-256-GCM, a 12-byte random IV per record and a 16-byte authentication
tag; 92.59\,\% statement coverage & Met \\
NFR-02 & I & \texttt{bcryptjs} at 10 rounds, 8-character minimum enforced at
registration; no recovery path exists & Met \\
NFR-03 & I & Session ID in a cookie with \texttt{httpOnly: true},
\texttt{secure: env.nodeEnv === 'production'}, \texttt{sameSite: 'lax'};
revocable via \texttt{sessionStore.destroy} / \texttt{destroyAllForUser} &
Met \\
NFR-04 & I & Credentials masked except through the FR-12 reveal action;
credential used as a scheduling grouping key is reduced to a one-way hash first
& Met (pending confirmation) \\
NFR-05 & I & Request-rate limiting is applied to the public survey endpoint with
distinct limits for form retrieval and submission; this is the only
unauthenticated endpoint family & Met \\
NFR-06 & T & AES-256-GCM authentication tag makes modification of the embedded
\texttt{surveyId}, \texttt{cycleId} or \texttt{deadline} fail decode rather than
succeed with altered values; \texttt{decodeToken} returns \texttt{null} rather
than throwing & Met \\
NFR-07 & I, A & No identity column or user foreign key on
\texttt{survey\_response} (structural); the shared-token and
non-joinable-recipient properties were verified by inspection
(\Cref{subsec:inv2}) & Met \\
NFR-08 & I & Aggregate survey insight is generated only once the number of
responses reaches the configured minimum
(\texttt{SURVEY\_MIN\_ANONYMOUS\_RESPONSES}, default 5, floor 3); below it, no
insight is produced & Met \\
NFR-09 & I & No personal identifiers found in the AI payload construction:
inspection for author, assignee, login, email, issue-key and branch fields found
only aggregate scores, depersonalised conditions, the project name and
respondent count. \textbf{However, free-text survey answers are transmitted
verbatim and may contain names the system cannot strip} (gap G-03) & Partial \\
NFR-10 & I & A failing connector does not abort the remaining connectors; the
run is recorded as partial (\texttt{failedTools[]}) & Met \\
NFR-11 & I & BullMQ retry with increasing backoff on every queue & Met \\
NFR-12 & T & Deterministic job identifier per project and time slot makes a
repeated or retried tick a no-op & Met \\
NFR-13 & I, A & Lexical search remains available as the ordinary mode, but it is
not used as failover: an unconfigured or failed Pinecone deep search returns
HTTP 503. The current tests do not assert failover. Survey question and insight
generation also have no non-AI substitute (\Cref{subsec:inv3}) & Partial \\
NFR-14 & I & Compensating logic on multi-table operations, required because the
Supabase REST client exposes no cross-table transaction (C-04) & Unverified \\
NFR-15 & T & The bounded worker pool limits parallel connector calls to four;
automated concurrency tests verify the bound & Met \\
NFR-16 & A & Projects sharing a credential staggered by
\texttt{SCHEDULED\_SYNC\_PAT\_SPACING\_MINUTES} (default 10) within
\texttt{SCHEDULED\_SYNC\_MAX\_SPREAD\_MINUTES} (default 180); GitHub connector
performs a pre-flight \texttt{rateLimit.remaining} check. Investigation INV-4
derives the headroom but lacks the per-connector call count & Partial \\
NFR-17 & I & Dashboard reads are served from stored snapshots; no dashboard
route calls an external tool & Met \\
NFR-18 & --- & A usability claim about the client's workflow that cannot be
established by inspection & Unverified \\
NFR-19 & D & Per-tool progress visible during a running synchronisation via SSE
& Met \\
NFR-20 & D, I & Any dimension score expands into its contributing signals with
values, weights and contributions & Met \\
NFR-21 & A & A connector registry maps each tool name to a factory behind a
shared interface; the risk engine consumes a typed metrics bundle rather than
provider-specific data. Six connectors coexist without changing the scoring
subsystem & Met \\
NFR-22 & T & Per-strategy and scoring-arithmetic unit tests; scoring rules
maintained in dedicated internal documentation & Met \\
NFR-23 & T & Type check, lint, build and tests are a required \texttt{ci-ok}
check before merge; the type check, lint and test commands were re-run
successfully for this evaluation & Met \\
NFR-24 & I & All configuration supplied through environment variables; no
credential or endpoint literal in source & Met \\
NFR-25 & I & Container configuration uses a Node~20 Alpine image and defines
Redis, API and worker services; the worker runs from the same application image
with an overridden command & Met \\
NFR-26 & I & Scheduling configured by named time zone
(\texttt{SCHEDULED\_SYNC\_TZ}, default UTC), independent of the host clock; a
migration converted naive timestamps to \texttt{timestamptz} & Met \\
\end{longtable}

\subsection{Evaluation Against Project Objectives}
\label{subsec:eval-objectives}

\begin{longtable}{
  >{\raggedright\arraybackslash}p{0.06\textwidth}
  >{\raggedright\arraybackslash}p{0.62\textwidth}
  >{\raggedright\arraybackslash}p{0.17\textwidth}
}
\caption{Achievement of project objectives}
\label{tab:objective-achievement}\\
\toprule
\textbf{ID} & \textbf{Basis of assessment} & \textbf{Outcome} \\
\midrule
\endfirsthead
\toprule
\textbf{ID} & \textbf{Basis of assessment} & \textbf{Outcome} \\
\midrule
\endhead
\bottomrule
\endlastfoot
O-01 & FR-13, FR-14, FR-24, FR-25 met: all four tool categories ingested and
presented in one view & \textbf{Achieved} \\
\addlinespace
O-02 & FR-20--FR-22 met; \Cref{subsec:inv1} establishes that the comparability
claim survives partial tool coverage, and that the rejected alternative would
have made tool coverage the dominant term in the score & \textbf{Achieved} \\
\addlinespace
O-03 & FR-26 and NFR-20 met: every score decomposes into named signals with
measured values, applied weights and individual contributions.
\Cref{subsec:inv1} records that this explainability is what ruled out the
imputation alternatives, since neither can state a breakdown in terms of the
project's own metrics & \textbf{Achieved} \\
\addlinespace
O-04 & NFR-07 and NFR-08 met: the response schema prevents direct identity
storage (INV-2 route A1), and the shared per-cycle token and non-joinable
recipient table close every inspected linkage route (\Cref{subsec:inv2}).
Respondent anonymity is enforced by the data model rather than promised by
policy & \textbf{Achieved} \\
\addlinespace
O-05 & FR-37--FR-38 and FR-40--FR-41 are met, while FR-39 is partial: actions
are recorded, retrievable lexically and optionally reranked through Pinecone,
but deep retrieval fails with 503 during provider unavailability. Whether
interventions \emph{demonstrably worked} cannot be assessed --- that requires a
deployment period longer than the project had, since effectiveness ratings only
become meaningful weeks after an action & \textbf{Partially achieved as a
mechanism} \\
\addlinespace
O-06 & FR-18 and NFR-12 met: scheduled synchronisation accumulates history
without user action, idempotently & \textbf{Achieved} \\
\addlinespace
O-07 & NFR-21 met: six connectors across four categories coexist without
modification to the scoring logic --- two more than the four in scope &
\textbf{Achieved, exceeded} \\
\end{longtable}

\subsection{Summary of Conformance}
\label{subsec:conformance-summary}

Of 41 functional requirements, \textbf{38 are met}, one is met pending a
confirmation the team must make (FR-31), and two are partial (FR-33, rotation
fairness unconfirmed; and FR-39, no deep-search failover). Of 26 non-functional
requirements, \textbf{21 are met}, three are partial (NFR-09, NFR-13, NFR-16)
and two are unverified by any admissible method (NFR-14, NFR-18). Of seven
objectives, five are achieved, one is achieved and exceeded (O-07), and one
(O-05) is partially achieved.

\subsection{Reproducing the Measurements}
\label{subsec:reproduction}

Every repository-derived quantitative figure in this chapter can be regenerated
using the commands in \Cref{tab:reproduction}.

\begin{longtable}{
  >{\raggedright\arraybackslash}p{0.38\textwidth}
  >{\raggedright\arraybackslash}p{0.52\textwidth}
}
\caption{Commands reproducing the measurements}
\label{tab:reproduction}\\
\toprule
\textbf{Figure} & \textbf{Command} \\
\midrule
\endfirsthead
\toprule
\textbf{Figure} & \textbf{Command} \\
\midrule
\endhead
\bottomrule
\endlastfoot
Type check (\Cref{tab:measurements}) & \texttt{cd backend \&\& npm run
typecheck} \\
Lint findings and warning counts & \texttt{cd backend \&\& npm run lint} \\
Vitest file/test counts and pass state & \texttt{cd backend \&\& npx vitest
run} \\
Coverage, overall and per-area (\Cref{tab:measurements,tab:coverage-by-area}) &
\texttt{cd backend \&\& npm run test:coverage} \\
Action-search suite (INV-3) & \texttt{cd backend \&\& npm run test:actions} \\
Test file and case counts & \texttt{find backend/apps backend/libs
backend/tests -name "*.test.ts" \textbar{} wc -l} \\
\end{longtable}

\Cref{subsec:inv1} cites \nolinkurl{backend/libs/risk-engines/scoring.test.ts}
(assertions at lines 48, 56, 68, 80 and 89) and the weight vector in
\nolinkurl{backend/libs/risk-engines/risks/security/security.strategy.ts}. The
worked example is arithmetic and can be checked by hand against those weights.

\section{Deviations and Design Revisions}
\label{sec:deviations}

Each deviation is given with its symptom, the technical analysis of its cause,
and the design revision made --- or the reason none was made.

\subsection{Transactional Email: Gmail SMTP Replaced by the Brevo HTTP API}
\label{subsec:deviation-email}

\paragraph{Symptom.} Transactional email --- verification codes, password
resets, invitations and survey links --- was initially delivered through
Nodemailer over Gmail SMTP and proved unreliable at the volume a multi-project
portfolio produces. A single survey dispatch sends one message per eligible
developer per project, and the scheduled distribution concentrates those sends
into a short window.

\paragraph{Cause.} A consumer mail account applies per-account sending limits
and reputation-based filtering that are not designed for machine-generated
transactional mail. Compounding this, SMTP introduces a stateful connection and
authentication handshake onto a path that must be stateless and cheaply
retryable from a background worker.

\paragraph{Revision.} Delivery was moved to Brevo's HTTP API. The email service
now uses a single idempotent HTTPS request instead of an SMTP connection,
allowing the existing queue retry policy (NFR-11) to handle transient provider
unavailability consistently. When no Brevo credential is configured, the service
records the intended message without sending it, so local development needs no
mail credential. No requirement changed; the mechanism satisfying it did.

\subsection{Removal of Manual Guidance from Survey Question Generation}
\label{subsec:deviation-guidance}

\paragraph{Symptom.} An earlier design allowed an administrator to inject
free-text guidance to steer AI question generation. In use, guidance text
propagated into the wording of the generated questions.

\paragraph{Cause.} This is a leading-question problem rather than a
prompt-engineering limitation, and could not have been addressed by better
prompting. A survey whose questions are shaped by a manager's stated concern
measures that concern rather than the team's state, and the resulting insight is
then presented to the same manager as an independent reading of team health. The
mechanism therefore undermined the objectivity that the risk-engine-driven
design exists to provide, and indirectly the candour on which the whole survey
subsystem depends.

\paragraph{Revision.} The override path was removed entirely. Questions now
derive strictly from computed risk scores, their trend and raw connector
metrics, as FR-27 specifies. Removing rather than constraining the feature was
deliberate: a constrained override would have had to be policed by judgement on
every use, whereas an absent one cannot be misused. The
\texttt{custom\_guidance} column survives on the \texttt{survey} table as
dormant schema and should be dropped in a future migration.

\subsection{Project Tracking State Moved to Persisted Storage}
\label{subsec:deviation-tracking}

\paragraph{Symptom.} The portfolio ``tracked''/bookmark state, initially held in
browser-local storage, did not survive a change of device or browser, so a
manager's curated portfolio view was silently lost.

\paragraph{Cause.} The state is a property of the user's relationship to a
project, not of the browser session, and had been placed at the wrong
architectural layer.

\paragraph{Revision.} \texttt{is\_tracked} was added as a persisted column,
exposed through the health feed and an administrator toggle endpoint, and the
portfolio filter and bookmark control were wired to read and write it.

\subsection{Scope Exceeded: Connectors and Notification Channels Beyond
Specification}
\label{subsec:deviation-connectors}

\Cref{ch:problem} scopes the work to one connector per category --- GitHub,
Jira, SonarCloud and GitHub Actions --- and scopes survey broadcast to Slack and
Discord. The delivered system contains \textbf{six} connectors (adding GitLab
under VCS and Linear under project management) and \textbf{three} notification
channels (adding Telegram), all behind the same interfaces.

This is an extension beyond scope rather than a shortfall, and it is positive
evidence for NFR-21 and objective O-07: additional providers were absorbed
without any change to the scoring logic. It is declared here because
\Cref{ch:problem}'s scope statement predates these additions.

\subsection{Acknowledged Limitation: Score Confidence Is Not Expressed}
\label{subsec:deviation-confidence}

\Cref{subsec:inv1} establishes that renormalisation removes the bias
zero-imputation would introduce, but it cannot recover the information an
absent signal would have carried: a score computed from two signals is a weaker
estimate than one computed from six, and the interface presents both
identically. Team health is most exposed, having only four signals, so one
absent input removes a quarter of the evidence behind the score.

No requirement demands otherwise, so this was recorded as a limitation rather
than non-conformance: a manager comparing two projects in the portfolio view
cannot see that one score rests on materially less evidence than the other,
which is precisely the comparison O-02 exists to support. The natural revision
is to surface the count or weight-fraction of present signals alongside each
score --- the \texttt{weights} array returned by
\texttt{renormalizedWeightedScore()} and the signal catalogue already carry all
the data required, so this is a presentation change rather than a model change.
It is recorded as future work in \Cref{ch:conclusion}.

\subsection{Deployment Topology}
\label{subsec:deviation-deployment}

\Cref{ch:research}'s feasibility study records that free hosting tiers covered
the API and frontend but not the always-on worker, which holds long-lived BullMQ
connections and therefore cannot run as a serverless function; it concluded that
the worker ``remains unhosted''. That conclusion was superseded during
development, and the system is now hosted and operational (C-08).

In the deployed configuration, requests under the API path are forwarded from
the Vercel-hosted client to the Render-hosted backend, while the worker also
runs on Render's always-on tier. The container build produces one application
image for both backend processes, with the worker selected through a command
override.
